\section{A diagonal argument on the external wires}\label{sec:diagonal}

Suppose that $k$ is the largest dimension of a spectral flat over a
finite field $F$. We shall construct two tensors on the same external wire
ends whose difference has price strictly below that of the through
identity $I$. We shall then compute this difference entry by entry and
obtain $I$ itself.

Each cell of a grid contains copies of a universal weighted RAM
computation, which we call a \emph{tick}. A program $a$, the controller,
specifies how to attach the wire ends to the ticks. Let $Z(a)$ be the
tensor computed by the exact grid. Theorem~\ref{thm:collapse} supplies
a smaller network $X$ for a group of ticks, with a price gap on the
error. Replacing each such group and its borrowed wires by $X$ gives a
second tensor $\widehat Z(a)$. Lemma~\ref{lem:lens} carries the local gap
to the external wire ends for every choice of $a$.

The program $a$ is chosen last. With $X$ fixed, an entry of
$\widehat Z(a)$ has a weighted evaluator that takes the controller code
as data. The saving in the size of $X$ leaves enough time for the
\emph{exact} ticks, acting together, to run this evaluator. We use one
branch of $a$ for that computation and one to transmit $I$. Their sum
will give $Z(a)=\widehat Z(a)+I$.

\subsection{The exact grid and its accelerated evaluator}

Put $V=2^w$, where $w$ tends to infinity through positive integers.
Use a word RAM with $64w$-bit words and addresses. Fix an instruction set and a universal interpreter. Each instruction on
boundedly many words has a uniformly constructible bit circuit of size
$w^{O(1)}$. Boolean operations, shifts, comparisons and ordinary word
arithmetic with fixed overflow conventions suffice. Memory starts at zero
apart from the input and advice, which are loaded one pair of bits per word.
A program can be supplied in a readable prefix of the input. The universal machine includes loading, instruction
fetch and decoding in its step count. Each fixed program fits in $w^2$
bits for large $w$. A weighted program may guess one bit in a step.
Its value is the sum in $F$ of its accepting terminal weights, with
rejection on invalid instructions or timeout. Simulating a capped machine
with a fresh address space costs at most a fixed power of $w$ per step,
for example by storing its visited addresses in a binary trie.

The parameters must meet two budgets. Each tick has $P$ pins on
which to receive its part of the input. Its time budget is $m$.
We spread the external wire pairs over enough grid cells to meet the pin
budget, and use the saving in Theorem~\ref{thm:collapse} to meet the
time budget.

Fix a sufficiently large integer saving exponent $A$. Its lower bound
will depend only on the polynomial overheads of the RAM simulation,
entry evaluation and trace distribution. Round $m\asymp V/w^{c_0}$ to a power of
two, where $c_0$ is large enough that the universal weighted RAM at cap
$m$ with $P\le m$ pin bits (including program) has a trivalent network
of size $\le V$. In the trace construction below, each of $O(m)$ rows
requires $w^{O(1)}$ gates, including those used to check memory accesses. The exponent $c_0$ is independent of $A$. For
Theorem~\ref{thm:collapse}, choose uniform envelopes $Q=V^{o(1)}$
on its replication factor $q$ and $D_*=V^{1-o(1)}\le V$ below its
pointwise gap. Before applying the theorem choose a pin length
$P=V^{1-o(1)}$ with
\[
         P w^c=o(V)\quad\text{for every fixed }c,\qquad
         QP=o(D_*).
\]
For example take $P=\lfloor D_*/(Q2^{\sqrt w})\rfloor$. This also
satisfies the theorem's input-pin limit, whatever fixed power of $w$
that limit uses. To expose the data, sum over the $2^{qP}$
assignments to the pins, placing their indicators at the common data
site. The star through that site supplies a summation diagonal of price
at most $2k$, so Equation~\eqref{eq:price-ordinary-sum} reduces the gap
by at most $2kqP$.
Thus a network $X$ approximates the product of $q$ universal ticks,
all pins exposed, with gap at least $D=D_*/2$. Pad by independent through
identities to give $X$ exactly
\[
                        B=\lceil qV/w^A\rceil
\]
catalyst pairs in each direction. This leaves its gap unchanged and
its size is $O_k(B)$, counting all ports. Over the fixed finite field
a list of its gates, weights and wiring fits in $B w^2$ bits, with
padding for all sufficiently large $w$.

Use the $k$-dimensional grid of Lemma~\ref{lem:lens} with $g$ cells
along each axis. Put $q$ tick slots per cell, for $T=g^kq$ slots.
Take $K=2g^{k-1}B$ external pairs in each direction, $E=kK$ altogether.
Choose $g=V^{o(1)}$ tending to infinity fast enough that
$Bw^c/g=o(P)$ for every fixed $c$. Choose $R\ge8g$, $R=V^{o(1)}$,
with $Eg^2/R=o(D)$. Such choices exist because $B/P$ and $Eg^2/D$
are at most $V^{o(1)}$ at their respective steps. One can dominate
each ratio by a factor $2^{\sqrt w}$. Then $Rg^k=o(D)$. Finally
choose a selector length $\ell=V^{1-o(1)}$ with
\[
                    R\ell=o(D),\qquad \ell=o(P),
\]
for example by diminishing $\min(P,D/R)$ by $2^{\sqrt w}$ and
rounding down. Each length or cap used by a program occupies $O(w)$
bits. The scale $R$ is used only for the estimate.

For a $w^2$-bit string $a$ define $Z(a)$ and $\widehat Z(a)$ as sums over
$\ell$-bit selectors, with external indices on the $E$ pairs. For a
selector run the deterministic controller encoded by $a$ for at most
$g^kB$ word steps. It can read $a$, the public lengths, the selector
and the description of $X$. Its settings are independent of the external
indices. Reject an invalid program, timeout or an illegal setting.
The controller may also reject explicitly. A rejected selector contributes
zero to both tensors. A legal setting obeys the following rules.
\begin{enumerate}
\item Each external pair is assigned a grid line in its own direction.
It may be touched at one cell on the line, feeding its incoming and
outgoing rail bits to two specified pins of a tick there. Pins used are
distinct. Every remaining pin has a prescribed constant bit.
\item Each line carries between $B$ and $3B$ pairs. At each cell of
that line at least $B$ are untouched there. The first $B$ such pairs
in label order serve as the cell's catalyst passages in this direction.
\end{enumerate}
In the exact tensor $Z(a)$ put the product of the $q$ ticks together
with their through identity at each cell. Untouched passages transmit
their bits by equality. A touch supplies two bits to the tick. Its
two sides can take independent values. In $\widehat Z(a)$ substitute
$X$ for the cell product and its through identity. This substitution
satisfies Lemma~\ref{lem:lens} separately for every legal term, even
when the catalyst passages vary along a line. At the sites $\pm R e_j$
of any spectral $k$-flat its estimate gives
\begin{equation}\label{eq:diagonal-price}
 H\bigl(Z(a)-\widehat Z(a)\bigr)
 \le 2RE-D+O_k(Eg^2/R)+2kR(\ell+g^k)<2RE=H(I).
\end{equation}

There is a fixed weighted evaluator for an entry of $\widehat Z(a)$
given $a$, the parameters, the $X$ description and the $2E$ external
bits. It guesses one of the $\ell$-bit selectors in the definition of
$\widehat Z(a)$, simulates and checks the controller,
and contracts the resulting network by guessing internal wire bits
and multiplying local entries. One can bucket pairs by line, check
pin assignments by direct addressing, and generate the rail passages
in line order. Each of the $g^k$ cells uses one listed copy of $X$, each of the $E$
rails has $g$ passages, and there are $TP$ pins in the settings.
Thus the path time is $(g^kB+gE+TP)w^{O(1)}$.
Here $gE=2kg^kB$ and $TP=o(g^kB)$. Take the public evaluator cap
$U\asymp g^kB w^{c_1}$ a power of two. The exponent $c_1$ can be
chosen independently of $A$, since the evaluator and its simulation
use the same algorithms for all public lengths. Its entire input and
advice, including its fixed program, have length $o(U)$, and
$\log_2 U\sim w$.

\subsection{Distributing the evaluator's trace}\label{app:trace}

To run this evaluator on the exact grid, we must split its work among
ticks. Both bits of an external pair must go to a single tick, since a
rail can supply them together at one cell. The controller must decide
where to attach the rails before it sees their bit values. We therefore
need a decomposition whose assignment of input bits depends only on a
short string. All running times for weighted routines below are
worst-case times along one path.

\begin{lemma}[Distribution with paired inputs]\label{lem:trace}
Suppose $U$ is a power of two with $\log_2 U\sim w$. Consider a capped
weighted computation of time $U$ with an input-advice string of at most
$U$ bits, divided into consecutive pairs with at most a final singleton.
For each integer $s\ge1$ with $s=o(w)$, the value on input-advice bits
$u$ has an exact expansion
\[
       \sum_h\sum_{t\in\{0,1\}^{L(h)}} \prod_j f_{h,j}\bigl(u|_{J_{h,j}},t\bigr).
\]
Here $h$ ranges over admitted compression histories of $s$ blocks, $t$
lists the bits on the final wires, and $j$ ranges over final compressor
positions. For each $h$, the sets $J_{h,j}$ partition the input-advice
positions, keeping each pair in one set.

For all sufficiently large $w$, each history has at most $V^{0.6}$ factors.
A canonical history and its numerical descriptors use at most $V^{0.6}$
bits. Checking the history and constructing the descriptors takes
deterministic time at most $V^{0.6}$, independently of the values of the
input bits and shared bits. The descriptors specify the input pairs as
lists of intervals and the shared string length $L(h)\le 2^{-s}U w^{O(1)}$.

Each $J_{h,j}$ has size at most $V^{0.6}$. A weighted subroutine evaluates
$f_{h,j}$ from the bits at these positions, with indexed access to $t$
and the descriptors. Its running time on each path is at most $W_j$, where the
descriptors give integer charges satisfying
\[
                  \sum_j W_j\le U w^{c_2},
                  \qquad \max_j W_j\le V^{0.6}.
\]
The work exponent $c_2$ depends only on the trace simulation and the
fixed compressor over $F$. It is fixed before choosing the exponent
$A$ in Theorem~\ref{thm:collapse}.
\end{lemma}

\begin{proof}
First we represent the computation by rows of local tensors. We then
group rows into a wire network to which Theorem~\ref{thm:compressor}
applies. At fixed final wire bits, each remaining position will yield
one factor in the expansion.

The RAM trace has a trivalent bit network with $O(U)$ rows, rounded
to a power of two, and $w^{O(1)}$ gates and incident bits per row.
Initialize two consecutive input bits in separate tracks of the same
row. Split each execution step into boundedly many microsteps making at
most one memory probe apiece, counting instruction fetches. As in the
sorting-based memory check for RAM traces \cite{TinyRAM}, guess reads,
check register transitions, and sort initialization and probe records
by address and time. Adjacent records enforce each read value, starting
with zero at an uninitialized address. Use deterministic local circuits,
canonical padding for unused fields and a unique continuation on halt.
Every accepting path has exactly one valid trace completion and its field
weight is supplied at the last row. In Batcher's bitonic sorting
network~\cite{Batcher}, each exchange layer uses fixed XOR partners.
Successive-row connections have XOR offsets that are strings of low
ones. The offset list has polynomial length in $w$, and the local gates,
partners and slot numbers are computable from row and track identifiers
in $w^{O(1)}$ steps. At any smaller cap $m$ of logarithm $O(w)$ the same construction
uses $O(m)w^{O(1)}$ gates, as required for the universal tick.

Put $S=2^{\lfloor w/2\rfloor}$ and group consecutive rows in buckets
of $S$. A bucket contains $V^{1/2+o(1)}$ gates, and there are
$V^{1/2+o(1)}$ buckets. This choice leaves room under the $V^{0.6}$
bounds for both computing a local factor and describing the entire
partition. The high row bits form a binary vector, or lane. At each lane put a
cyclic cursor through its coordinate positions. From a cursor site use
two arcs to the next cursor site, staying on the lane or flipping the
current coordinate. Repeat arcs to obtain the fixed power-of-two in-
and out-degree $d_0$ of Theorem~\ref{thm:compressor}. Base endpoints are
distinct and the displacement of each base arc is fixed. Place the bucket
tensor at cursor zero. Route each inter-bucket XOR offset through a
cursor cycle, flipping its indicated coordinates. Each step is a
permutation of the lanes. Give each offset slots numbered by exchange
type, low source row bits and bit within the message. Relays equate
the occupied slots and pin unused ones to zero. The prefix of the
offset already traversed determines the original lane for a relay.
Pad arc widths uniformly, reserving slots for every such combination of
offset, exchange type, low source row bits and bit number. Their ranges
give the capacities in $w^{O(1)}$ work per position.
It follows that the number of positions $P_0$ and maximum arc width $W_0$
satisfy
\[
         P_0\le (U/S)w^{O(1)},\qquad W_0\le Sw^{O(1)}.
\]
To evaluate a bucket or relay on given arguments, generate its rows,
guess the local internal bits and check the wires. This takes
$Sw^{O(1)}$ steps per path.

Proposition~\ref{lem:ports} supplies port graphs over $F$, of fixed
degree and arbitrarily large power-of-two order, in the loopless symmetric
form required by Theorem~\ref{thm:compressor}. Fix the resulting compressor
identities over $F$.
The number of lane coordinates is $\log_2(U/S)+O(1)\sim w/2$, enough
for $s$ blocks, each consuming a fixed number. For the theorem's mass
$1/(d_0P_0)$ per arc the initial square root of energy is at most
$d_0P_0W_0\le Uw^{O(1)}$. The final total arc width is at most
$2^{-s}(d_0P_0W_0+O(P_0))$. Concatenate the final arc strings into $t$, using the same slice of
$t$ at the two ends of each arc. At fixed $t$ the contraction is the
product over final positions. An admitted history $h$ sends each whole
old factor to one new position at every substep. Thus the input bits
of a given row go to one $J_{h,j}$, with both bits of an external pair
still together.

For the numerical descriptors use the scans in Lemma~\ref{lem:descriptors}.
They generate bounded lists of arc endpoints, widths, offsets, ties,
slices, preimages and map records at each position. Motif admission uses
binary widths as in Theorem~\ref{thm:conversion}, and layout comparisons use
sums of squares with fixed rational coefficients after removal of the
common scale. The number of blocks is $o(w)$, with a bounded number of
substeps each, and all widths and counts stay polynomial in $U$.
They therefore take $O(w)$ bits. Use fixed-length indices or lengths
determined by earlier choices, and canonical padding. Following the
unique destination of each initial bucket gives its final factor's
interval list without inspecting any input value. There are at most
$P_0$ bucket intervals. The combined time and storage are
$P_0w^{O(1)}=V^{1/2+o(1)}$.

Evaluate each factor down its tree of constituent tensors, guessing
the summed fields and multiplying local entries. Charge the initial
bucket costs above. At each substep charge a fixed polynomial in $w$
times the number of positions, the total field lengths and the sum
of $1+M_0$ over local map occurrences, where $M_0$ is the maximum
demand width of the corresponding motif. In particular the map entries
of Theorem~\ref{thm:conversion} can be evaluated by scanning the $O(1+M_0)$
fixed-size selectors, checking copies and totals, and multiplying
the selected base entries with running offsets on the bounded list of
fields. The resulting work is $(1+M_0)w^{O(1)}$ per map. Each field
or maximum demand uses boundedly many incident widths, and the
constituent records go to unique final positions. The block-boundary
energy estimate in Theorem~\ref{thm:compressor}, and its two substep estimates
\eqref{eq:compress-shrink}, \eqref{eq:compress-refresh} within each block, bound total width at every substep by
$Uw^{O(1)}$. Summing the charges gives the claimed total, with room
for indexing, subroutine cleanup and sequential multiplication of
several factors. Per factor there are at most $C^s$ constituent objects
of width at most $C^s(W_0+1)$ for a fixed $C$, by the same theorem.
It follows that the work and the input-bit count at each position are
$V^{1/2+o(1)}\le V^{0.6}$ for large $w$. The bounds were obtained by
counting fields, widths and constituent records. These are numerical
descriptors, so the controller can compute each charge before the input
bits arrive.
\end{proof}


\subsection{Programming the exact computation}

All three overhead exponents $c_0,c_1,c_2$ were fixed independently
of $A$. Choose $A>c_0+c_1+c_2$, so the mean work charge fits within a tick's time budget.
\[
       \frac{Uw^{c_2}}{Tm}
       =O\!\left(w^{c_0+c_1+c_2-A}\right)=o(1).
\]
Apply
Lemma~\ref{lem:trace} to the evaluator, initializing the two bits
of every external pair together, and using fixed padding between fields
where needed. Take $s=o(w)$ large enough that
$2^{-s}U w^{O(1)}=o(\ell)$. This is possible since
$\log_2(U/\ell)=o(w)$. The history also has length $o(\ell)$.
Reserve a fixed-length public input field for $s$, so choosing its
value preserves all earlier cap and input-length bounds. No program
needs to compute which $s$ suffices.

The controller uses exactly one selector for an identity branch. It
distributes pairs equally over lines in each direction, touches none
and gives every tick a program returning one. Every other accepted selector encodes a history $h$ and final wire
bits $t$ from Lemma~\ref{lem:trace}. A branch bit, descriptor-specified
lengths and checked padding give exactly one selector per pair $(h,t)$.
The evaluator's own guess of a selector for $\widehat Z(a)$ is part of
the computation expanded by the lemma.

For each such selector distribute the factors of Lemma~\ref{lem:trace}
among the $T$ ticks. The resources of a factor are its work charge, its
number of external pairs in each of the $k$ directions, and its number
of other input-advice bits. Each coordinate is at most $V^{0.6}$.
For each coordinate separate zero
and group positive integers in intervals $[b^i,b^{i+1})$, $b=21/20$.
Group factors by the tuple of their intervals and distribute each group
round-robin over the ticks. There are $O_k(w^{k+2})$ groups. In a
single group the number of factors at a tick differs from the group's
average by at most one. If the total mean load per tick in a coordinate
is $L$, its load at every tick is therefore between
\[
             L/b-O_k(V^{0.6}w^{k+2})
           \quad\text{and}\quad
             bL+O_k(V^{0.6}w^{k+2}).
\]
Indeed each positive coordinate in a group lies between its lower
endpoint and $b$ times that endpoint. Summing the per-group estimate
gives the bounds. Computing the classes takes $w^{O_k(1)}$ work per
factor, for example by precomputing the integer endpoints with exact
arithmetic. Factor counts, descriptors and assignment can all be
processed in time at most $V^{0.8}$ for large $w$.

Assign each external pair to the direction-appropriate line through
its tick's cell, with its unique touch at that tick. In each direction
the mean load per tick is $K/T=2B/(gq)=V^{1-o(1)}$, dominating the
additive allocation error. A line has $gq$ ticks, hence carries
between $(2/b-o(1))B$ and $(2b+o(1))B$ pairs. Its $q$ ticks at one
cell touch $O(B/g)$ of them. These bounds ensure at least $B$ untouched
passages at that cell and a total load between $B$ and $3B$ on the line.

Each tick receives its fixed program, index, parameters, selector
and the input bits assigned to its factors. The last are provided by
touches for external pairs and by controller-supplied constants for all
others, in the interval orders from the descriptors. In particular, the descriptors identify the portion of the evaluator
program, the description of $X$, and the controller string $a$ needed
at this tick. The controller sets the corresponding pins. The assigned bit
count at a tick is
\[
       O_k\bigl((E+Bw^2)/T+V^{0.6}w^{k+2}\bigr)=o(P).
\]
Here the partition in Lemma~\ref{lem:trace} divides the actual
input-advice string, of length $O_k(E+Bw^2)$. Each tick also receives
the shared selector, using $o(P)$ pins. Pad all remaining pins. To emit
the settings it suffices to describe the line and unique touch of
each pair and give the pin assignments. Constructing the descriptors
and allocation, checking the selector and writing these settings costs
at most
\[
             V^{0.8}+(TP+E+B)w^{O(1)}=o(g^kB).
\]
Here $g$ dominates any fixed power of $w$, and
$Pw^c=o(V)$ for every fixed $c$. One line assignment and one touch assignment per pair specify all of
its passages.

The receiving tick regenerates the descriptors and the allocation from
its selector and parameters, operating on the interval lists. By using
their lengths and offsets it locates its own inputs without scanning the
other ticks' input values. Its setup time is at most
$V^{0.8}+Pw^{O(1)}=o(m)$. It evaluates its factors successively on
weighted paths, multiplying their terminal weights. By the allocation
bound its work is at most $bUw^{c_2}/T+O_k(V^{0.6}w^{k+2})=o(m)$
in addition to setup. The charged evaluations include local addressing
and reuse of scratch memory.

Choose $a$ to encode precisely this controller. The tick network and
its approximation $X$ were chosen for the universal machine before $a$
was fixed. The accelerated evaluator accepts $a$ as data, and the
controller reads the bytes of $a$ from its own inputs when specifying
the exact ticks. The evaluator interprets any $w^2$-bit program up to its public cap.
The layout of its RAM trace depends only on the universal interpreter
and the public lengths. Allocation and local evaluation are therefore
specified before the values of the advice bits, including the program
bits, are known. Once the field and compressor are fixed, the routines
are fixed programs with the prescribed padding. The controller can
supply the ticks with the required bytes of these programs and of its
own input $a$.

Fix the external bits. On a simulation selector, equality along each
rail forces the bits on each side of its unique touch, with exactly
one assignment between that touch and each external endpoint. For the selected $(h,t)$, each tick multiplies the factors assigned to
it, summing over the internal paths of their weighted evaluations. Thus
the product over ticks is the corresponding product in
Lemma~\ref{lem:trace}. Summing over $(h,t)$ gives the specified entry
of $\widehat Z(a)$, and the unique remaining selector contributes the
entry of $I$. This holds at every entry over $F$, so
\[
           Z(a)-\widehat Z(a)=I,\qquad
           H\bigl(Z(a)-\widehat Z(a)\bigr)=2RE .
\]
This contradicts the strict inequality in
Equation~\eqref{eq:diagonal-price}. Spectral flats over each finite
field therefore have unbounded dimension. Restricting to coordinate
subspaces gives flats of every finite dimension, and the field transfer
in Lemma~\ref{lem:transfer} proves Theorem~\ref{thm:flats} over every
field.

\begin{proof}[Completion of the proof of Theorem~\ref{thm:main}]
The finite-dimensional flats just obtained supply the hypothesis of
Theorem~\ref{thm:hilbert}. Its compatible prices on any real Hilbert
space assign each bit wire its length, respect congruences of labeled
placements, and are homogeneous under dilation for nonzero tensors.
Any nonempty finite set $S$ of sites lies in a finite-dimensional affine
subspace. On this subspace Theorem~\ref{thm:flatlaw} gives the price of a
common $q$-valued label on $S$ as $(\log_2 q)W(S)$, where $W(S)$ is the
Gaussian mean width normalized to give a segment its length. When
$S$ has positive diameter, the spectral point on $S$ takes the value
$2$ on $\Delta_S(2)$. Its base-two logarithm therefore has multiplier
$H(\Delta_S(2))=W(S)$.
\end{proof}
